% ============================================================
%  Independent Implementation and Evaluation of Dynamic Rank
%  Distribution for Parameter-Efficient Fine-Tuning of RoBERTa
%
%  Compile with pdfLaTeX on Overleaf.
%  All figure files are optional: a framed placeholder is
%  inserted automatically when the PNG is absent.
% ============================================================
\documentclass[conference]{IEEEtran}

% ----- Standard IEEE-compatible packages --------------------
\usepackage{graphicx}
\usepackage{amsmath}
\usepackage{amssymb}
\usepackage{booktabs}
\usepackage{array}
\usepackage{cite}
\usepackage{url}
\usepackage{algorithm}
\usepackage{algpseudocode}

% Robust figure placeholder: show framed box if PNG absent
\usepackage{ifthen}
\usepackage{fancybox}

% Helper: insert image or a framed placeholder of the same
% nominal size when the file does not exist.
% Usage: \figplaceholder{filename}{width}{height}{caption text}
\newcommand{\figplaceholder}[4]{%
  \IfFileExists{#1}{%
    \includegraphics[width=#2]{#1}%
  }{%
    \framebox[#2]{\parbox[c][#3][c]{0.85\linewidth}{\centering
      \small\textbf{[Figure placeholder]}\\[2pt]
      Filename: \texttt{#1}\\[2pt]
      #4
    }}%
  }%
}

\begin{document}

% ============================================================
%  Title and Author
% ============================================================
\title{Independent Implementation and Evaluation of Dynamic Rank
Distribution for Parameter-Efficient Fine-Tuning of RoBERTa}

\author{%
  \IEEEauthorblockN{[Author Name]}
  \IEEEauthorblockA{%
    [Department Name]\\
    [Institution Name]\\
    [City, Country]\\
    [email@institution.edu]
  }
}

\maketitle

% ============================================================
%  Abstract
% ============================================================
\begin{abstract}
Fine-tuning large pretrained language models in their entirety
imposes significant memory and computational demands, motivating
a class of approaches termed parameter-efficient fine-tuning
(PEFT). This report presents an independent educational
implementation of dynamic rank distribution, a PEFT method
introduced by Chen et al.\ \cite{chen2024dora}, which augments
a frozen pretrained model with trainable low-rank components
and adaptively removes those components that contribute least
to the learned transformation. The implementation adapts
RoBERTa-base \cite{liu2019roberta} for binary sentiment
classification on the Stanford Sentiment Treebank version 2
(SST-2) dataset \cite{wang2018glue}. Seventy-two linear
transformations across all twelve encoder layers are replaced
by rank-3 adaptive modules, giving an initial rank budget of
216 components across the full network. A cubic pruning
schedule, informed by component-level importance scores
accumulated via exponential moving averaging, reduces this
budget to a target of 144 components over the first half of
training. A diversity-enhancing mutual-information-inspired
regularization term is applied throughout to promote
non-degenerate factor distributions. After sixty training
epochs, the implementation achieves 99.74\% accuracy on the
full SST-2 training set and 93.69\% accuracy on the
validation set, with the final rank budget of 144 matching
the prescribed target exactly, representing a 33.33\%
reduction from the initial allocation.
\end{abstract}

% ============================================================
%  Index Terms
% ============================================================
\begin{IEEEkeywords}
Parameter-efficient fine-tuning, Transformer, RoBERTa,
low-rank adaptation, dynamic rank distribution, structured
pruning, sentiment analysis, regularization, natural language
processing.
\end{IEEEkeywords}

% ============================================================
%  I. Introduction
% ============================================================
\section{Introduction}

The proliferation of large pretrained language models such as
BERT \cite{devlin2019bert}, RoBERTa \cite{liu2019roberta}, and
GPT-family architectures has transformed natural language
processing. These models encode broad linguistic knowledge in
hundreds of millions of parameters acquired during
unsupervised pretraining on large text corpora. When adapted
to a downstream task, conventional fine-tuning updates every
parameter in the network, a strategy that demands substantial
GPU memory, extended training time, and separate storage of a
full model copy per task. These costs become prohibitive as
model scale increases.

Parameter-efficient fine-tuning (PEFT) addresses this by
keeping pretrained weights frozen and introducing a small
number of additional trainable parameters. Low-rank adaptation
(LoRA) \cite{hu2022lora} is among the most widely adopted PEFT
strategies: it inserts a pair of low-rank matrices alongside
each target linear transformation, allowing the weight update
to be learned at a fraction of the original parameter count.
A key limitation shared by naive low-rank methods is the
assumption that the same rank allocation is appropriate for
every layer and every component within a layer. Empirical
analysis suggests that the information content of different
linear transformations varies substantially, making uniform
rank allocation suboptimal.

The DoRA method \cite{chen2024dora} addresses this limitation
by decomposing each low-rank update into individually gated
rank-1 components, measuring the relative importance of each
component throughout training, and progressively removing
those judged least informative. The final model retains a
globally optimized, non-uniform rank budget that concentrates
capacity where it is most needed.

This report documents an independent educational reproduction
of these ideas. The work does not claim to replicate every
experimental result of the original publication; rather, it
aims to implement the core algorithmic components faithfully,
verify them through unit testing, and evaluate the resulting
system on a well-understood benchmark. The principal
contributions of this project are:
\begin{itemize}
  \item An independent PyTorch implementation of adaptive
    rank-1 component decomposition, component importance
    estimation, exponential moving-average score tracking, a
    cubic rank-budget schedule, global pruning with
    deterministic tie-breaking, permanent final masking, and
    variance-based regularization.
  \item Integration of the resulting adapter modules into a
    frozen RoBERTa-base backbone for SST-2 classification.
  \item A training and evaluation pipeline including
    checkpointing and automatic resume support.
  \item Experimental results demonstrating that the dynamic
    rank pruning process reaches the prescribed final budget
    of 144 total active components, achieving 93.69\%
    validation accuracy.
\end{itemize}

The remainder of the report is structured as follows.
Section~\ref{sec:background} provides the theoretical
foundations. Section~\ref{sec:math} presents the mathematical
formulation. Section~\ref{sec:arch} describes the system
architecture and implementation. Section~\ref{sec:exp} covers
the experimental setup. Section~\ref{sec:results} reports the
experimental outcomes. Section~\ref{sec:discussion} discusses
the findings. Section~\ref{sec:limitations} identifies
limitations and future directions. Section~\ref{sec:conclusion}
concludes the report.

% ============================================================
%  II. Background and Theoretical Foundations
% ============================================================
\section{Background and Theoretical Foundations}
\label{sec:background}

\subsection{Artificial Intelligence and Machine Learning}
Artificial intelligence (AI) refers to computational systems
designed to perform tasks that would otherwise require human
judgment. Machine learning (ML), a central subdiscipline of
AI, constructs models that improve their performance on a
task by learning statistical regularities from data, rather
than through explicit rule specification \cite{goodfellow2016dl}.

\subsection{Supervised Learning and Binary Classification}
Supervised learning trains a model on input-output pairs,
minimizing a loss function that measures disagreement between
predicted and true outputs. Sentiment analysis on SST-2 is a
binary classification task: given a movie-review sentence, the
model must predict one of two classes, positive or negative.
The model is evaluated by the proportion of correctly
classified examples, referred to as classification accuracy.

\subsection{Transfer Learning and Pretrained Language Models}
Training a large model from random initialization on a
small-scale supervised task is data-inefficient and often
produces poor generalization. Transfer learning leverages a
model pretrained on a large general-purpose corpus and
subsequently adapts it to the target task. Transformer-based
language models such as BERT \cite{devlin2019bert} and
RoBERTa \cite{liu2019roberta} are pretrained using
self-supervised objectives on massive text collections,
acquiring representations that transfer effectively to diverse
downstream tasks.

\subsection{The Transformer Architecture}
The Transformer \cite{vaswani2017attention} is a neural
network architecture built from stacked blocks, each
containing a multi-head self-attention sublayer and a
position-wise feed-forward sublayer. Self-attention computes
three linear projections of the input sequence, referred to as
queries ($\mathbf{Q}$), keys ($\mathbf{K}$), and values
($\mathbf{V}$), and produces a weighted combination of value
representations guided by the dot-product similarity between
queries and keys. The feed-forward sublayer applies two
successive linear transformations separated by a non-linear
activation function.

In a 12-layer RoBERTa-base encoder, each block contains
linear projections for the query, key, value, and output of
the attention mechanism, as well as two projections in the
feed-forward sublayer, yielding six linear transformations per
block. Input text is first tokenized into sub-word tokens,
padded or truncated to a uniform sequence length, and
accompanied by an attention mask that indicates which token
positions carry genuine content.

\subsection{Forward Propagation, Loss, and Optimization}
During a forward pass, the model transforms a batched input
through all encoder layers to produce class logits. The
cross-entropy loss measures the disagreement between the
predicted class distribution and the true one-hot labels.
Gradients of the loss with respect to all trainable parameters
are computed via backpropagation and used to update the
parameters. The AdamW optimizer \cite{loshchilov2019adamw}
combines adaptive moment estimation with decoupled weight
decay, and is widely used for Transformer fine-tuning.

\subsection{Parameter-Efficient Fine-Tuning}
PEFT methods freeze the pretrained weights and introduce a
small set of trainable parameters. This reduces both memory
consumption (no need to maintain optimizer states for all
backbone parameters) and storage overhead (task-specific
parameters are orders of magnitude smaller than the full
backbone). LoRA \cite{hu2022lora} achieves this by adding a
product of two low-rank matrices alongside each target linear
layer; the pretrained weight is unchanged during fine-tuning.

\subsection{Low-Rank Matrix Factorization and Rank-1 Components}
A matrix of rank $r$ can be expressed as the sum of $r$
rank-1 outer products. In the context of PEFT, restricting the
trainable weight update to a low-rank subspace constrains the
degrees of freedom available to the optimizer and acts as an
implicit regularizer. When the update is decomposed into
individual rank-1 components, each component can be assigned a
trainable scalar gate that controls its contribution
independently, enabling fine-grained selection of the
informative subspace.

\subsection{Importance-Based Pruning and Exponential Moving
Averages}
Structured pruning removes entire units (neurons, channels, or
rank components) from a model based on an importance criterion,
typically the magnitude of their contribution to the output.
To reduce sensitivity to instantaneous fluctuations, the
importance score is accumulated across training steps using an
exponential moving average (EMA), which smooths the signal by
weighting recent observations more heavily while retaining
historical context.

\subsection{Cubic Pruning Schedules}
A pruning schedule determines when and by how much the rank
budget should be reduced during training. A cubic schedule
concentrates the majority of the reduction early in the
pruning phase (because the cubic function has a steep initial
slope) while tapering off toward the end, giving the model
time to adapt to the reduced capacity before pruning is
finalized.

\subsection{Variance-Based Regularization and Structured
Sparsity}
Variance-based regularization penalizes low-variance factor
matrices, discouraging the collapse of multiple rank-1
components onto a single dominant direction. This promotes
diversity among the learned components, which can be
interpreted as a form of implicit orthogonality enforcement.
Structured sparsity refers to setting entire units, rather
than individual weights, to zero; zeroing a scalar gate $c_i$
renders an entire rank-1 component inactive while leaving its
underlying parameter tensors intact.

\subsection{Training and Validation Accuracy, Overfitting}
Training accuracy is the fraction of training examples
classified correctly by the model evaluated at its final
state. Validation accuracy measures generalization to unseen
examples. A substantial gap between the two figures suggests
overfitting: the model has memorized training-set patterns
that do not transfer to the validation set. Checkpointing
saves the complete training state at defined intervals,
enabling recovery from hardware or software failures without
discarding the accumulated training progress.

% ============================================================
%  III. Mathematical Formulation
% ============================================================
\section{Mathematical Formulation}
\label{sec:math}

\subsection{Frozen Pretrained Weights and Adaptive Updates}
Let $W_0 \in \mathbb{R}^{d_{\text{out}} \times d_{\text{in}}}$
denote the pretrained weight matrix of a linear layer. During
adapter training, $W_0$ is kept frozen and a trainable low-rank
residual $\Delta W$ is added:
\begin{equation}
  W_{\text{eff}} = W_0 + \Delta W.
  \label{eq:eff_weight}
\end{equation}
The forward computation therefore produces
$W_{\text{eff}}\,\mathbf{x}$ for any input vector
$\mathbf{x}$, where the pretrained knowledge is preserved in
$W_0$ and task-specific adaptation is carried by $\Delta W$.

\subsection{Rank-1 Component Decomposition}
The weight update is expressed as an explicit sum of
rank-1 outer products, each controlled by an independent
trainable scalar gate:
\begin{equation}
  \Delta W
    = \frac{\alpha}{r}\sum_{i=1}^{r} c_i\,
      \mathbf{b}_i \mathbf{a}_i^{\top},
  \label{eq:decomp}
\end{equation}
where $\mathbf{a}_i \in \mathbb{R}^{d_{\text{in}}}$ is the
$i$-th row of the input-side factor matrix
$A \in \mathbb{R}^{r \times d_{\text{in}}}$,
$\mathbf{b}_i \in \mathbb{R}^{d_{\text{out}}}$ is the $i$-th
column of the output-side factor matrix
$B \in \mathbb{R}^{d_{\text{out}} \times r}$,
$c_i \in \mathbb{R}$ is the $i$-th scalar gate, $r$ is the
maximum rank, and $\alpha / r$ is an adapter scaling factor.
This can be written compactly as
$\Delta W = (\alpha/r)\,B\,\mathrm{diag}(\mathbf{c})\,A$,
which matches the implementation in \texttt{layer.py}
(line~194: \texttt{low\_rank * self.scale}, where
\texttt{scale = alpha / rank}).

\subsection{Adapter Scaling}
The scalar $\alpha / r$ normalizes the magnitude of the update
relative to the rank, following the convention established in
LoRA \cite{hu2022lora}. In the project configuration,
$\alpha = 2.0$ and $r = 3$, giving a scaling factor of
$2/3 \approx 0.667$. This value is computed once at
initialization and stored as \texttt{self.scale}.

\subsection{Forward Computation}
For an input tensor $X$ of shape
$[\text{batch}, \text{seq\_len}, d_{\text{in}}]$, the full
forward pass computes:
\begin{equation}
  Y = W_0 X^{\top}
    + \frac{\alpha}{r}\,\bigl((X\,A^{\top}) \odot
      \mathbf{c}^{\top}\bigr)\,B^{\top},
  \label{eq:forward}
\end{equation}
where $\odot$ denotes element-wise multiplication applied
across the rank dimension. Setting $c_i = 0$ eliminates the
contribution of component $i$ without modifying $A$ or $B$.

\subsection{Component Importance}
The relative importance of component $i$ within a layer is
measured by comparing its contribution to the total weight
update using Frobenius norms:
\begin{equation}
  s_i = \frac{\|\Delta W_i\|_F}
             {\left\|\sum_{j=1}^{r}\Delta W_j\right\|_F + \varepsilon},
  \label{eq:importance}
\end{equation}
where $\Delta W_i = (\alpha/r)\,c_i\,\mathbf{b}_i
\mathbf{a}_i^{\top}$ is the rank-1 weight matrix of component
$i$, and $\varepsilon = 10^{-12}$ is a small constant added
for numerical stability to prevent division by zero when all
gates are near zero. This formula is implemented in
\texttt{importance.py}.

\subsection{Exponential Moving Average}
Raw importance scores computed from the current parameter
state may be noisy. To obtain a stable estimate, an EMA is
maintained for each component across training steps:
\begin{equation}
  \hat{s}_i(t) = \beta\,\hat{s}_i(t-1)
               + (1 - \beta)\,s_i(t),
  \label{eq:ema}
\end{equation}
where $\beta = 0.9$ is the decay coefficient. This biases
the estimate toward the historical average, reducing the
influence of transient parameter configurations on the
pruning decisions.

\subsection{Cubic Rank-Budget Schedule}
Pruning proceeds only within a defined window of training:
from step $t_{\text{start}}$ to step $t_{\text{end}}$. The
target average rank per layer at step $t$ is:
\begin{equation}
  b(t) = \begin{cases}
    b_0 & t \leq t_{\text{start}}, \\[4pt]
    b_0 - (b_0 - b_T)\,\xi^3 & t_{\text{start}} < t < t_{\text{end}}, \\[4pt]
    b_T & t \geq t_{\text{end}},
  \end{cases}
  \label{eq:schedule}
\end{equation}
where $b_0$ is the initial average rank per layer, $b_T$ is
the target final average rank, and
$\xi = (t - t_{\text{start}}) / (t_{\text{end}} -
t_{\text{start}}) \in [0, 1]$ is the normalized progress
within the pruning window. The cubic exponent causes the rank
to decrease rapidly early in the window and more gradually
near the end, mirroring the typical convergence behavior
described in \cite{chen2024dora}. In this project,
$b_0 = 3$, $b_T = 2$, $t_{\text{start}} = 0.15 \times T$,
and $t_{\text{end}} = 0.50 \times T$, where $T = 63{,}180$
total training steps.

\subsection{Global Pruning}
At each pruning checkpoint, the EMA scores from all adaptive
layers are concatenated into a single vector of length equal
to the total number of rank components across the network.
The target total number of active components $K_t$ is
computed as $K_t = \lfloor b(t) \times L \rceil$, where $L$
is the number of adaptive layers. The bottom
$(K_{\max} - K_t)$ components are identified using
\texttt{torch.argsort} with \texttt{stable=True}, which
provides deterministic tie-breaking when multiple components
share identical scores (such as at initialization when all
gates are zero). The corresponding scalar gates are set to
zero, temporarily suppressing those components without
modifying the underlying $A$ and $B$ tensors.

During the active pruning phase, suppressed components remain
trainable: gradient descent can restore a gate from zero. At
the final pruning checkpoint $t_{\text{end}}$, a boolean
permanent mask is recorded based on which gates are zero,
and the method \texttt{enforce\_final\_mask()} is called
immediately after each subsequent optimizer step to re-apply
this mask, preventing any gradient-driven recovery.

\subsection{DEM Regularization}
The diversity-enhancing mutual-information (DEM) regularization
term \cite{chen2024dora} penalizes degenerate factor matrices
by measuring the unbiased sample variance of each row of $A$
and each column of $B$:
\begin{equation}
  L_{\text{DEM}} =
    \frac{1}{N_c}\sum_{\ell}\left(
      \sum_{i=1}^{r}\mathrm{Var}(\mathbf{a}_i^{(\ell)})
      + \sum_{i=1}^{r}\mathrm{Var}(\mathbf{b}_i^{(\ell)})
    \right),
  \label{eq:dem}
\end{equation}
where the outer sum is taken over all adaptive layers $\ell$,
$\mathrm{Var}(\mathbf{a}_i)$ is the unbiased variance of the
$i$-th row of $A$ (computed over the $d_{\text{in}}$
elements), $\mathrm{Var}(\mathbf{b}_i)$ is the unbiased
variance of the $i$-th column of $B$ (computed over the
$d_{\text{out}}$ elements), and $N_c = 2 r L$ is the total
count of variance terms across all layers and both factors.
The combined training objective is:
\begin{equation}
  L_{\text{total}} = L_{\text{task}} + \eta\,L_{\text{DEM}},
  \label{eq:total_loss}
\end{equation}
where $L_{\text{task}}$ is the cross-entropy classification
loss and $\eta = 0.5$ is the regularization coefficient.

% ============================================================
%  IV. System Architecture and Implementation
% ============================================================
\section{System Architecture and Implementation}
\label{sec:arch}

\subsection{Overall Architecture}

The project is organized as a Python package rooted at
\texttt{src/dora/} with the following module responsibilities:
\begin{itemize}
  \item \texttt{layer.py} — defines \texttt{AdaptiveRankLinear},
    the core module implementing Eq.~\eqref{eq:decomp}
    and Eq.~\eqref{eq:forward}.
  \item \texttt{importance.py} — implements
    \texttt{component\_importance()} computing
    Eq.~\eqref{eq:importance}.
  \item \texttt{scheduler.py} — implements
    \texttt{CubicBudgetScheduler} computing
    Eq.~\eqref{eq:schedule}.
  \item \texttt{pruner.py} — implements
    \texttt{DynamicRankPruner}, which manages EMA state
    (Eq.~\eqref{eq:ema}), global pruning, and the final mask.
  \item \texttt{regularization.py} — implements
    \texttt{dem\_regularization()} computing
    Eq.~\eqref{eq:dem} and \texttt{dem\_loss()} computing
    Eq.~\eqref{eq:total_loss}.
  \item \texttt{transformer.py} — implements
    \texttt{freeze\_model()} and
    \texttt{replace\_linear\_layers()}.
  \item \texttt{examples/hf\_roberta\_sst2.py} — training
    script integrating all modules for the RoBERTa/SST-2
    experiment.
\end{itemize}

\begin{figure}[t]
  \centering
  \figplaceholder{figures/overall_workflow.png}{0.95\linewidth}{5cm}{%
    Flowchart: SST-2 dataset $\rightarrow$ tokenizer $\rightarrow$
    frozen RoBERTa $\rightarrow$ adaptive low-rank layers
    $\rightarrow$ task loss + DEM regularization
    $\rightarrow$ optimizer $\rightarrow$ importance
    calculation $\rightarrow$ EMA $\rightarrow$ cubic
    pruning $\rightarrow$ final rank distribution
    $\rightarrow$ evaluation.}
  \caption{Overall project workflow from raw dataset to
    final evaluation.}
  \label{fig:workflow}
\end{figure}

\subsection{Adaptive Linear Layer}

Each target \texttt{nn.Linear} module is replaced by an
\texttt{AdaptiveRankLinear} instance that stores three
trainable parameter groups: matrix $A \in
\mathbb{R}^{r \times d_{\text{in}}}$, matrix $B \in
\mathbb{R}^{d_{\text{out}} \times r}$, and scalar gate vector
$\mathbf{c} \in \mathbb{R}^r$. Both $A$ and $B$ are
initialized using Kaiming uniform initialization; $\mathbf{c}$
is initialized to zero, so the adapter contributes no update
at the start of training. The frozen base weights are stored
in a separate \texttt{nn.Linear} submodule with
\texttt{requires\_grad=False}.

\begin{figure}[t]
  \centering
  \figplaceholder{figures/adaptive_layer.png}{0.95\linewidth}{4.5cm}{%
    Diagram: frozen $W_0$; trainable $A$, $B$, $c_i$;
    scaled low-rank residual; combined effective weight.}
  \caption{Structure of a single \texttt{AdaptiveRankLinear}
    module.}
  \label{fig:adaptive_layer}
\end{figure}

\subsection{RoBERTa Integration and Layer Targeting}

Following \texttt{freeze\_model()}, which sets
\texttt{requires\_grad=False} on every existing parameter, the
function \texttt{replace\_linear\_layers()} traverses the
model's module hierarchy by exact name and substitutes the six
target linear projections per encoder block. For a 12-layer
RoBERTa-base model, the six targets per block are:
\texttt{attention.self.query},
\texttt{attention.self.key},
\texttt{attention.self.value},
\texttt{attention.output.dense},
\texttt{intermediate.dense}, and
\texttt{output.dense}.
This yields $12 \times 6 = 72$ adaptive layers.

\begin{figure}[t]
  \centering
  \figplaceholder{figures/roberta_adaptation.png}{0.95\linewidth}{5cm}{%
    Diagram of 12-layer RoBERTa encoder with six adapted
    linear projections per block: Q, K, V, attention
    output, intermediate dense, output dense.}
  \caption{RoBERTa-base adaptation: six linear
    transformations per encoder block replaced by adaptive
    rank-1 component layers.}
  \label{fig:roberta}
\end{figure}

\subsection{Importance Estimation and EMA Tracking}

After each optimizer step, \texttt{DynamicRankPruner.step()}
calls \texttt{update\_importance()}, which computes
$s_i(t)$ for every layer using Eq.~\eqref{eq:importance}
and updates the stored EMA score using
Eq.~\eqref{eq:ema}. All EMA scores are maintained on CPU
as float32 tensors.

\subsection{Global Pruning and Budget Scheduling}

At each pruning checkpoint (every 10 steps within the
pruning window, plus forcibly at $t_{\text{end}}$), the
pruner collects EMA scores from all 72 layers into a single
flat tensor. It computes the required number of active
components as $K_t = \lfloor b(t) \times 72 \rceil$ and
determines the number to remove as $K_{\max} - K_t$. Using
\texttt{torch.argsort} with \texttt{stable=True}, it
selects the lowest-scoring components and sets their
corresponding $c_i$ values to zero.

\begin{figure}[t]
  \centering
  \figplaceholder{figures/dynamic_pruning.png}{0.95\linewidth}{5cm}{%
    Flowchart: importance calculation $\rightarrow$ EMA
    smoothing $\rightarrow$ global score collection
    $\rightarrow$ budget schedule $\rightarrow$ argsort
    selection $\rightarrow$ temporary gate zeroing
    $\rightarrow$ final mask at $t_{\text{end}}$.}
  \caption{Dynamic rank pruning mechanism.}
  \label{fig:pruning}
\end{figure}

\subsection{Final Pruning Mask}

At step $t_{\text{end}}$, the boolean tensor recording which
gates are zero is stored as the final mask per layer.
Subsequently, \texttt{enforce\_final\_mask()} is called
immediately after every \texttt{optimizer.step()}, zeroing
out any gate that the optimizer might have updated away from
zero. This ensures that the final topology is stable
throughout the remainder of training. Critically, the
underlying $A$ and $B$ parameter tensors of pruned components
are not deleted; only their gates are frozen at zero. Physical
tensor removal would require post-training restructuring of
the weight matrices, which is not currently implemented.

\subsection{DEM Regularization in Practice}

The function \texttt{dem\_regularization()} iterates over all
\texttt{AdaptiveRankLinear} modules, computes
\texttt{torch.var(layer.A, dim=1, unbiased=True)} (one value
per row of $A$, i.e., per rank component) and
\texttt{torch.var(layer.B, dim=0, unbiased=True)} (one value
per column of $B$), sums all terms, and normalizes by
$N_c = 2rL = 2 \times 3 \times 72 = 432$. The function is
differentiable with respect to $A$ and $B$, ensuring that
gradients flow through the regularization term.

\subsection{Training Pipeline and Checkpointing}

The training loop iterates over mini-batches; for each step
it performs a forward pass, computes the combined loss
(Eq.~\eqref{eq:total_loss}), runs backpropagation, applies
the AdamW optimizer, calls \texttt{pruner.step()}, and
enforces the final mask. A checkpoint is saved to
Google Drive after every completed epoch, recording the
model state dictionary, optimizer state dictionary, pruner
state (EMA scores, final mask, \texttt{pruning\_finished}
flag), the completed epoch number, the global step count,
and the Python, NumPy, and PyTorch RNG states. On restart,
the script automatically detects the checkpoint and resumes
training from the following epoch.

\subsection{Training Pseudocode}

Algorithm~\ref{alg:train} summarizes the complete training
and pruning procedure as implemented.

\begin{algorithm}[t]
\caption{Dynamic rank adaptation training procedure}
\label{alg:train}
\begin{algorithmic}[1]
\Require Frozen backbone $\Theta_0$, adaptive layers
  $\{A_\ell, B_\ell, \mathbf{c}_\ell\}_{\ell=1}^{L}$,
  initial rank $r$, final rank $r_T$, epochs $E$, batch
  size $B$, learning rate $\eta_0$, DEM coefficient $\eta$,
  EMA decay $\beta$, pruning start/end fractions
\Ensure Trained adapters with final rank budget $r_T L$
\State Initialize all $c_{\ell,i} \leftarrow 0$;
  $\hat{s}_{\ell,i} \leftarrow 0$
\State $T \leftarrow$ total training steps; $t \leftarrow 0$
\For{epoch $= 1$ to $E$}
  \For{each mini-batch $(X, y)$}
    \State $\hat{y} \leftarrow$ \textsc{Forward}($X$)
      \Comment{Eq.~\eqref{eq:forward}}
    \State $L_{\text{task}} \leftarrow
      \textsc{CrossEntropy}(\hat{y}, y)$
    \State $L_{\text{DEM}} \leftarrow
      \textsc{DemReg}(\{A_\ell, B_\ell\})$
      \Comment{Eq.~\eqref{eq:dem}}
    \State $L \leftarrow L_{\text{task}} + \eta\,L_{\text{DEM}}$
    \State \textsc{Backward}($L$)
    \State \textsc{AdamW.Step}()
    \For{each adaptive layer $\ell$}
      \State $s_{\ell,i}(t) \leftarrow
        \textsc{Importance}(\ell)$
        \Comment{Eq.~\eqref{eq:importance}}
      \State $\hat{s}_{\ell,i}(t) \leftarrow
        \beta\,\hat{s}_{\ell,i}(t{-}1)
        + (1{-}\beta)\,s_{\ell,i}(t)$
        \Comment{Eq.~\eqref{eq:ema}}
    \EndFor
    \If{\textsc{ShouldPrune}($t$)}
      \State Compute $K_t$ from Eq.~\eqref{eq:schedule}
      \State Sort all $\hat{s}_{\ell,i}$ globally
      \State Set $c_{\ell,i} \leftarrow 0$ for
        bottom $(K_{\max} - K_t)$ components
      \If{$t = t_{\text{end}}$}
        \State Record permanent mask; set
          \texttt{pruning\_finished}$= \mathrm{True}$
      \EndIf
    \EndIf
    \If{\texttt{pruning\_finished}}
      \State \textsc{EnforceFinalMask}()
    \EndIf
    \State $t \leftarrow t + 1$
  \EndFor
  \State \textsc{SaveCheckpoint}(epoch, $t$, state)
\EndFor
\end{algorithmic}
\end{algorithm}

% ============================================================
%  V. Dataset and Experimental Setup
% ============================================================
\section{Dataset and Experimental Setup}
\label{sec:exp}

\subsection{Dataset}
The Stanford Sentiment Treebank version 2 (SST-2)
\cite{socher2013recursive}, distributed as part of the GLUE
benchmark \cite{wang2018glue}, is a binary sentiment
classification dataset derived from movie reviews. Each
example consists of a single sentence annotated as either
positive or negative. The dataset was accessed via the
Hugging Face \texttt{datasets} library using the identifier
\texttt{nyu-mll/glue, sst2}. The full training split was
used for training (TRAIN\_SUBSET\_SIZE was set to
\texttt{None}, meaning no subset selection was applied),
and the official validation split was used for evaluation.
Exact split sizes should be verified against the dataset
documentation; they are not independently reproduced here.

\subsection{Preprocessing}
The RoBERTa-base tokenizer was applied to each sentence with
truncation to a maximum length of 128 sub-word tokens.
Batches were assembled using \texttt{DataCollatorWithPadding},
which pads sequences to the length of the longest sequence
within each batch.

\subsection{Model Configuration}
RoBERTa-base \cite{liu2019roberta} was loaded from the
Hugging Face model hub with a sequence classification head
initialized for two output classes. All pretrained parameters
were frozen. Adaptive modules with rank $r = 3$ and scaling
parameter $\alpha = 2.0$ were inserted into the six target
linear transformations of each of the 12 encoder blocks, for
a total of 72 adaptive layers and a maximum total rank of
$72 \times 3 = 216$. Adapter dropout was set to 0.0.

\subsection{Training Hyperparameters}

Table~\ref{tab:hyperparams} lists the full training
configuration.

\begin{table}[t]
  \centering
  \caption{Training configuration}
  \label{tab:hyperparams}
  \begin{tabular}{lc}
    \toprule
    \textbf{Hyperparameter} & \textbf{Value} \\
    \midrule
    Optimizer            & AdamW \\
    Learning rate        & $8 \times 10^{-4}$ \\
    Batch size           & 64 \\
    Number of epochs     & 60 \\
    Initial rank $r$     & 3 \\
    Target final rank $r_T$ & 2 \\
    Scaling $\alpha$     & 2.0 \\
    Adapter dropout      & 0.0 \\
    DEM coefficient $\eta$ & 0.5 \\
    EMA decay $\beta$    & 0.9 \\
    Pruning start fraction & 0.15 \\
    Pruning end fraction & 0.50 \\
    Pruning interval     & 10 steps \\
    Maximum sequence length & 128 \\
    \bottomrule
  \end{tabular}
\end{table}

\subsection{Pruning Schedule}
With $T = 63{,}180$ total steps (derived from the training
set size and batch size), pruning began at step
$t_{\text{start}} = \lfloor 0.15 \times 63180 \rfloor =
9{,}477$ and concluded at step
$t_{\text{end}} = \lfloor 0.50 \times 63180 \rfloor =
31{,}590$. The target average rank per layer decreased
from $b_0 = 3$ to $b_T = 2$ over this interval according to
Eq.~\eqref{eq:schedule}.

\begin{figure}[t]
  \centering
  \figplaceholder{figures/rank_schedule.png}{0.95\linewidth}{4cm}{%
    Schematic graph of total rank budget vs.\ training step.
    Horizontal: flat at 216 before $t_{\text{start}}$;
    cubic decay from 216 to 144 between $t_{\text{start}}$
    and $t_{\text{end}}$; flat at 144 thereafter.
    Note: schematic only; exact per-step values not recorded.}
  \caption{Schematic of the cubic rank-budget schedule.
    Exact step-by-step values were not recorded during
    training and would need to be re-derived from the
    scheduler for an exact reproduction.}
  \label{fig:schedule}
\end{figure}

\subsection{Hardware and Software}
Training was carried out on an NVIDIA Tesla T4 GPU via Google
Colab. The software stack comprised PyTorch (version to be
verified from the training environment), the Hugging Face
Transformers and Datasets libraries, and Python 3. Software
version numbers were not logged during training and should be
confirmed for full reproducibility.

\subsection{Evaluation Protocol}
After training concluded, the final model checkpoint was
loaded and the final pruning mask was enforced. The model was
then evaluated in inference mode (\texttt{model.eval()},
\texttt{torch.no\_grad()}) on the full training set and on the
official validation set. Accuracy is defined as the proportion
of examples for which the argmax of the output logits matches
the ground-truth label.

% ============================================================
%  VI. Experimental Results
% ============================================================
\section{Experimental Results}
\label{sec:results}

Table~\ref{tab:results} summarizes the experimental outcomes
recorded from the final checkpoint.

\begin{table}[t]
  \centering
  \caption{Final experimental results}
  \label{tab:results}
  \begin{tabular}{lc}
    \toprule
    \textbf{Metric} & \textbf{Value} \\
    \midrule
    Model                      & RoBERTa-base \\
    Dataset                    & SST-2 \\
    Training epochs            & 60 \\
    Global training steps      & 63,180 \\
    Initial total rank         & 216 \\
    Final total active rank    & 144 \\
    Target final total rank    & 144 \\
    Training accuracy          & 99.74\% \\
    Validation accuracy        & 93.69\% \\
    \bottomrule
  \end{tabular}
\end{table}

\subsection{Rank Reduction}
The pruning procedure reduced the total active rank from 216
to 144, achieving the prescribed target exactly. This
represents a reduction of 72 components, or 33.33\% of the
initial budget. Because pruning operates globally rather than
per-layer, the retained components are distributed
non-uniformly across the 72 adaptive layers: individual layers
may retain different numbers of active components depending on
their measured importance during training. The exact per-layer
rank distribution is illustrated schematically in
Fig.~\ref{fig:rank_dist}; precise values would require
inspection of the \texttt{pruner.summary()} output from the
saved checkpoint.

\begin{figure}[t]
  \centering
  \figplaceholder{figures/final_rank_distribution.png}{0.95\linewidth}{4cm}{%
    Bar chart: active rank (0, 1, 2, or 3) for each of the
    72 adaptive layers after final pruning. Values should be
    read from \texttt{pruner.summary()} of the saved
    checkpoint. Insert actual data before submission.}
  \caption{Final rank distribution across the 72 adaptive
    layers. Exact per-layer values must be read from the
    saved checkpoint and inserted here.}
  \label{fig:rank_dist}
\end{figure}

\subsection{Accuracy}
The model achieved 99.74\% accuracy when evaluated on the
full SST-2 training set and 93.69\% accuracy on the
validation set. Both figures were obtained by running the
evaluation mode (\texttt{--evaluate} flag) of the training
script against the epoch-60 checkpoint. No epoch-by-epoch
accuracy curves were recorded; only the final values are
available.

\begin{figure}[t]
  \centering
  \figplaceholder{figures/accuracy_comparison.png}{0.95\linewidth}{3.5cm}{%
    Bar chart with two bars: Training accuracy 99.74\%
    and Validation accuracy 93.69\%.}
  \caption{Final training and validation accuracy on SST-2.}
  \label{fig:accuracy}
\end{figure}

% ============================================================
%  VII. Discussion
% ============================================================
\section{Discussion}
\label{sec:discussion}

\subsection{Parameter Efficiency}
By restricting the trainable weight update to rank-3
components of selected linear layers, the implementation
drastically reduces the number of parameters updated during
fine-tuning relative to full fine-tuning of RoBERTa-base.
The exact trainable parameter count depends on the
dimensions of the replaced layers and would need to be
extracted from the model before stating a precise ratio.

\subsection{Effectiveness of Dynamic Rank Allocation}
The global pruning budget reached the prescribed target of
144 components exactly, confirming that the cubic schedule,
importance scoring, and EMA are working correctly. The
non-uniform distribution of retained ranks across layers is
an expected outcome of global selection: layers whose
components carry more informative gradients naturally retain
more capacity, while those that converge to near-zero
contributions are pruned more aggressively.

\subsection{Training and Validation Accuracy Gap}
The gap between training accuracy (99.74\%) and validation
accuracy (93.69\%) is approximately 6 percentage points,
which is consistent with a degree of overfitting. Several
factors may contribute to this: the learning rate of
$8 \times 10^{-4}$ is relatively high for Transformer
fine-tuning; no learning-rate schedule or warm-up was
applied (based on the configuration described); and no
dropout was applied in the adapter modules. Establishing the
precise cause would require additional ablation experiments
that were not conducted in this study.

It is important to note that a 93.69\% validation accuracy on
SST-2 is a competitive figure for a PEFT method, and is
broadly consistent with results reported in the PEFT
literature for similar adapter configurations. However, no
controlled baseline experiments (e.g., fixed-rank LoRA, full
fine-tuning) were performed in this project, so a direct
quantitative comparison cannot be made.

\subsection{Rank Reduction Versus Physical Compression}
The implementation represents rank reduction through gate
zeroing rather than actual tensor deletion. The $A$ and $B$
matrices of pruned components remain allocated in GPU memory
and are traversed during the forward pass. Achieving actual
memory and latency savings would require restructuring the
weight matrices to exclude the zero-gated components, a step
not currently implemented. The rank budget therefore measures
the conceptual sparsity of the adapter, not a directly
measurable reduction in inference cost.

\subsection{Reproducibility and Checkpointing}
The checkpoint system proved valuable during the training run:
the ability to resume from any completed epoch without
discarding the accumulated training state (including EMA
scores and the final pruning mask) is critical for long
multi-epoch runs on transient cloud infrastructure.

% ============================================================
%  VIII. Limitations and Future Work
% ============================================================
\section{Limitations and Future Work}
\label{sec:limitations}

\subsection{Limitations}
\begin{itemize}
  \item \textbf{Single dataset and configuration.} Results are
    reported only for SST-2 with one specific rank budget.
    Generalization to other tasks, languages, or model scales
    cannot be inferred.
  \item \textbf{Absence of controlled baselines.} No
    comparison against full fine-tuning, fixed-rank LoRA, or
    other PEFT methods was performed; relative performance
    claims are therefore unsubstantiated.
  \item \textbf{No statistical replication.} A single training
    run was conducted. Multiple random seeds are required to
    characterize variance in the results.
  \item \textbf{Small rank budget.} The target rank of 2 per
    layer is at the lower end of practical values; behavior
    at larger budgets is uncharacterized.
  \item \textbf{Training-validation accuracy gap.} The 6
    percentage-point gap suggests overfitting; the
    contributing hyperparameters were not ablated.
  \item \textbf{No inference efficiency measurements.} GPU
    memory usage, inference latency, and FLOPs were not
    measured; claims about deployment efficiency would be
    speculative.
  \item \textbf{Soft sparsity only.} Zeroed gates do not
    reduce physical memory footprint or computation unless
    followed by tensor restructuring.
  \item \textbf{No precision, recall, or F1 reporting.} Only
    accuracy was measured; a complete classification
    evaluation would include these metrics.
\end{itemize}

\subsection{Future Work}
\begin{itemize}
  \item Conduct baseline experiments using fixed-rank LoRA and
    full fine-tuning to establish quantitative comparisons.
  \item Explore multiple rank budgets (e.g., target ranks of
    1, 2, 4, 8) and measure accuracy against rank.
  \item Evaluate on additional GLUE benchmarks and other
    tasks (e.g., question answering, named entity
    recognition).
  \item Report precision, recall, F1-score, and confusion
    matrices.
  \item Measure GPU memory consumption and wall-clock
    inference latency before and after pruning.
  \item Conduct multi-seed experiments to characterize result
    variance.
  \item Implement physical tensor compaction of pruned
    components to realize actual memory and latency savings.
  \item Apply a learning-rate schedule and warm-up to
    potentially reduce the training-validation accuracy gap.
\end{itemize}

% ============================================================
%  IX. Conclusion
% ============================================================
\section{Conclusion}
\label{sec:conclusion}

This report has presented an independent educational
implementation of the dynamic rank distribution method
introduced in \cite{chen2024dora}, applied to parameter-
efficient fine-tuning of RoBERTa-base on the SST-2 binary
sentiment classification benchmark. The implementation covers
the full algorithmic pipeline: rank-1 component decomposition
with scalar gate control, Frobenius-norm importance scoring,
exponential moving-average score smoothing, a cubic rank-budget
schedule, global deterministic pruning with permanent final
masking, and variance-based DEM regularization.

After 60 training epochs totaling 63,180 gradient steps, the
dynamic pruning procedure successfully reduced the total
active rank from an initial budget of 216 to the prescribed
final budget of 144, a reduction of 33.33\%. The final
checkpoint achieves 99.74\% accuracy on the full SST-2
training set and 93.69\% on the validation set. The exact
per-layer rank distribution is non-uniform, reflecting the
globally optimized pruning decisions guided by component
importance throughout training.

The project demonstrates that dynamic, importance-guided rank
allocation can be implemented reliably from first principles
and integrated with a frozen large pretrained Transformer
backbone. Future work should extend the evaluation to multiple
datasets, baselines, and random seeds, and should implement
physical weight compaction to realize the inference efficiency
benefits implied by the reduced rank budget.

% ============================================================
%  References
% ============================================================
\bibliographystyle{IEEEtran}
\bibliography{references}

\end{document}
